\section{Exact verification of the two certificates}
\label{sec:verification}

We now give the finite calculations establishing
Propositions~\ref{prop:pair-certificate} and~\ref{prop:mixed-certificate}.
The complete program and all its data are printed in
Appendix~\ref{app:code}.  It specifies the integer
matrices, coordinate orders, and finite enumerations used below.  All
arithmetic that certifies an equality or inequality is integer or rational
arithmetic.  Modular arithmetic supplies candidate substitutions, whose
validity is then checked exactly.
Rational Gram certificates and exact positivity checks are established
methods; see Peyrl--Parrilo \cite[Section~3]{PeyrlParrilo2008}.
Our input matrices are specified integers, and the proof below certifies
the displayed identities directly from those data.

The calculation has three steps.  First, expand each expression in named
moments using the identities already proved.  Next, use a finite set of
true linear equations to substitute for some of these moments.  Finally,
check either that every coefficient vanishes, for the pair certificate, or
that the constant coefficient plus the absolute sum of the remaining
coefficients is negative, for the mixed certificate.  We first justify the
substitutions for arbitrary real moment values, then specify the equations and coefficient checks for each certificate.

\subsection{Exact certification of linear substitutions}

Both certificates use finite systems of linear equations in their moment
coordinates. The same substitution procedure applies to either system.
Include the constant moment as a coordinate, so that every moment equation is
homogeneous in the enlarged coordinate vector.  Clearing denominators
turns the equations into integer rows.  The lists called \texttt{I} in
the program specify a fixed priority order on the coordinates.  This
choice affects the form of the answer and, in the mixed calculation,
the size of its residual coefficient sum; it does not add any relation.

\begin{lemma}[Exact substitutions]
\label{lem:exact-elimination}
Let $W\in\Z^{r\times N}$ consist of selected moment equations after
their denominators have been cleared.  Suppose a set $J$ of $r$ columns gives a square
submatrix $W_J$ invertible modulo a prime.  Put
$F=\{1,\ldots,N\}\setminus J$.  If a rational matrix
$P\in\Q^{N\times N}$ satisfies
\[
 P_{F,F}=I,\qquad P_{*,J}=0,\qquad WP=0,
\]
then every real vector $\mathbf m$ satisfying $W\mathbf m=0$ obeys
$\mathbf m=P\mathbf m$.  Thus the rows of $P$ give valid substitutions
for all coordinates in terms of the coordinates indexed by $F$.
\end{lemma}

\begin{proof}
Invertibility modulo the prime implies that the integer
$\det W_J$ is nonzero.  In the coordinate order $J,F$, write
\[
 W=\begin{pmatrix}W_J&W_F\end{pmatrix},\qquad
 P=\begin{pmatrix}0&P_{J,F}\\0&I\end{pmatrix}.
\]
The exact equality $WP=0$ gives
$P_{J,F}=-W_J^{-1}W_F$.  Every solution of $W\mathbf m=0$ consequently
satisfies $\mathbf m_J=P_{J,F}\mathbf m_F$, while the free coordinates
of $P\mathbf m$ are already $\mathbf m_F$.  This proves the claim over
$\mathbb R$, regardless of whether the moment values are rational.
Equivalently, $P$ has rank $N-r$, its image equals $\ker W$, and its
identity rules make it the identity on that kernel.
\end{proof}

Here is how \texttt{eliminate} verifies the hypotheses of the lemma.
Its first modular elimination uses the prime $1000000007$.  The routine
\texttt{echelon} retains an original integer equation whenever that
equation supplies a new pivot.  These retained equations form $W$.
Each pivot is at the largest remaining coordinate of its row; normalized
elimination therefore proves that $W_J$ is invertible modulo the prime.
Back substitution processes pivot coordinates in ascending order.
When a smaller pivot is removed, its already reduced row contains only
that pivot and free coordinates.  The final rule for a pivot thus has
no other pivot on its right side.

The program then works modulo the further moduli listed in the source,
retaining matching pivot sets, combines coefficients by the Chinese
remainder algorithm, and uses \texttt{rat} to propose rational rules.
These steps preserve the absence of other pivot coordinates.  All
nonpivot coordinates receive explicit identity rules.  Finally, for
each retained integer equation the program substitutes every proposed
rule, clears all denominators, and checks that every integer
coefficient is zero.  This is the exact check $WP=0$ in
Lemma~\ref{lem:exact-elimination}.

Every actual moment vector satisfies the selected equations, so this
check suffices even if additional generated equations have not been
used.  In particular, the modular pivot count is not asserted to be the
full rational rank of the original system.  A prime that conceals an
additional independent equation merely leaves that true equation
unused.  The rational reconstruction and its size threshold are ways
to find a candidate $P$; the exact check and the lemma establish its
validity.

\subsection{Pair equations and the vanishing identity}

For an integer array $A$ with zero row and column sums, put
\[
 \operatorname{wt}(A)=\sum_i\ord(A_{i,*}).
\]
This also equals the sum of the column orders.  Arrays name the real
injective moments $m_H(A)$ of Section~\ref{sec:pair}; the empty array names
the constant moment $1$.  A formal linear expression in these moments is
stored as a finite list of rational coefficients indexed by arrays.
Coefficients at the same index are added, and zero coefficients are
discarded.  These are the operations \texttt{add} and \texttt{lin} in the
program.

The routine \texttt{canonpair} selects a representative using precisely
the symmetries of $m_H$: permutation of either axis, removal of zero
rows and columns, global negation, and transposition.  More explicitly,
after removing zero axes it considers $A,-A,A^{\mathsf T},-A^{\mathsf T}$,
retains orientations with at least as many rows as columns, orders
columns by their positive-part orders, and minimizes over the remaining
column permutations after lexicographically sorting the rows.  Each
operation preserves the moment by its definition and
Lemma~\ref{lem:contraction}.  The degree orders and lexicographic choices
serve only to select consistent names.

Define $U_0(A)$ recursively as a linear expression in representatives
whose nonzero row and column orders are all at least two.  If a row of
$A$, or of $A^{\mathsf T}$, has order one, apply
Lemma~\ref{lem:contraction} in that orientation:
\begin{equation}
 U_0(A)=-\frac{1}{7-p}\sum_{\ell\ne i}U_0(A_{i\to\ell}),
 \qquad p\le6.
 \label{eq:recursive-pair-reduction}
\end{equation}
Canonicalize each contracted array before continuing.  If no such row
exists on either axis, retain its moment as a single coordinate.  The
empty array is retained in the same way.  Every contraction reduces the
sum of the two axis lengths, so this recursion terminates.  This is
\texttt{evpair}; \texttt{expand\_degree1} supplies the unsigned sum of
contractions.  The program checks that every evaluated array has both
dimensions at most six and has zero sums on both axes.

The pair table is formed from formal seven-row equations.  The exact
enumeration restrictions are as follows:
\begin{enumerate}[label=(\roman*)]
 \item Choose a total weight $k\in\{7,8,9,10,11\}$ and a column count
 $q$ with $2\le q\le\min(4,\lfloor k/2\rfloor)$.
 \item Choose nondecreasing sequences of seven row orders
 $(r_0,\ldots,r_6)$ and $q$ column orders $(c_0,\ldots,c_{q-1})$.
 Both sequences sum to $k$, and
 \[
 1\le r_i\le\lfloor k/2\rfloor,
 \qquad 2\le c_j\le\lfloor k/2\rfloor.
 \]
 \item Enumerate integer $7\times q$ arrays with zero row and column
 sums and exactly these orders.  Rows of equal order may be put in
 lexicographic order, and arrays with the same canonical representative
 may be skipped.
\end{enumerate}
The routines \texttt{parts}, \texttt{comps}, and \texttt{allAB} implement
this enumeration.  A row is generated as the difference of disjointly
supported nonnegative compositions of its prescribed order, within the
remaining positive and negative column capacities.  The capacities
enforce the prescribed column orders.  The first row has order one:
seven positive integer row orders with sum at most eleven cannot all
be at least two, and the orders are sorted.

For each retained array, the seven-row boundary case of
Lemma~\ref{lem:contraction} gives the equation
\begin{equation}
 \sum_{\ell=1}^{6}U_0(A_{0\to\ell})=0.
 \label{eq:pair-table-equations}
\end{equation}
Here $0,\ldots,6$ index the formal rows, not coordinates of a charge.
Only the contracted arrays, which have at most six rows, are evaluated.
The checks in \texttt{pairtable} verify the balance, number of formal
rows, and order of the first row before forming this equation.
Skipping an array can only omit an available true equation; the proof
does not require a complete set of moment relations.

Apply Lemma~\ref{lem:exact-elimination} to these equations.  Write
$U(A)$ for the resulting expansion of $U_0(A)$ in the retained moment
coordinates.  For a product used here, with charge columns
$a_1,\ldots,a_n$ and $n\le4$, let
$A_{\mathcal S}$ be the array obtained by adding the columns in each
block of a set partition $\mathcal S$ of $\{1,\ldots,n\}$.  The final
pair expansion is
\begin{equation}
 \sum_{\substack{\mathcal S\text{ a set partition}\\|\mathcal S|\le6}}
 \frac{(6)_{|\mathcal S|}}{6^n}\,U(A_{\mathcal S}).
 \label{eq:verified-partition-expansion}
\end{equation}
This is exactly Lemma~\ref{lem:partitions}.  The routine
\texttt{partition} enumerates set partitions by numbering blocks in
order of first occurrence, and \texttt{rep\_pair} implements
\eqref{eq:verified-partition-expansion}.  The routine \texttt{apair}
first forms the array of charge columns.  It may remove a zero charge
before the expansion, since that removes the factor $g_H(0)=1$.
Deletion of zero coordinate rows is also legitimate: a uniform
injection on the retained positions has the required uniform marginal.

For the pair certificate, \texttt{check\_pair} forms the displayed
polynomial $L_H$ of Proposition~\ref{prop:pair-certificate}, with
coefficients $-1,2,1,1,1$.  It multiplies each monomial by the conjugate
of every monomial, expanding $3|L_H|^2$ through
\eqref{eq:verified-partition-expansion}.  For each of the two stated
charges $q$, it adds the expansion with charge columns
$(\alpha,q,-\alpha,-q)$.  After the exact pair substitutions, the
result has no nonzero coefficient, for either $q$.  This proves the
identity in Proposition~\ref{prop:pair-certificate}; all expressions
used in the check have been identified with their actual averaged
moment values above.

\subsection{Mixed keys and the finite mixed table}

We use the notation $\mathcal M(\mathbf s;\mathbf a)$ and $m(A)$ from
Section~\ref{sec:mixed}.  The latter is the average of $m_H(A)$ over
the ordered basis pairs.  A mixed key consists of a label tuple
$\mathbf s$ and an array $A$ whose columns are the charges.  The array
has at most six nonzero coordinate rows.  Its rows and columns sum to
zero.  The routine \texttt{cmix} chooses representatives by simultaneous
factor and label permutations, renaming labels by order of first
occurrence, sorting the common coordinate rows, global negation, and
removing zero factors and zero coordinate rows.  These are exactly the
symmetries justified in Lemma~\ref{lem:mixed-reductions}.  The separate
tag $(\varnothing,A)$ names the pair moment $m(A)$; the completely empty
key names $1$.

The routine \texttt{evmix} supplies the following recursive expansion
of a key into named moment coordinates:
\begin{enumerate}[label=(\roman*)]
 \item If a charge has order one or is a permutation of $\alpha$, the
 mixed moment is zero.
 \item When at most one label remains, expand the product by
 \eqref{eq:verified-partition-expansion}, now using the averaged pair
 moments $m(A)$, and attach the pair tag to its coordinates.
 If this single-label expression was zero by the preceding rule,
 also record its pair expansion as a zero equation.  These recorded
 equations are stored in \texttt{rows0}.
 \item Otherwise, if an order-two charge has a label occurring exactly
 once among the factors, let $l$ be the number of labels in use and
 replace that label in turn by each of the other $l-1$ labels.  The
 original moment is $-1/(7-l)$ times the sum of the resulting moments.
 Apply the same rules recursively to each term.
 \item If none of these rules applies, retain the key as an original
 mixed-moment coordinate.
\end{enumerate}
The zero rules and the singleton coefficient follow from
Lemma~\ref{lem:mixed-reductions}.  Each singleton replacement decreases
the number of distinct labels.  The single-label use of the pair table
is legitimate because the real average over ordered pairs includes
the required transpose symmetry, as proved in that lemma.

The finite enumeration of triples is given by \texttt{triples(11)}.
For precision, it generates balanced charges $a,b,c\in\Z^6$ with
\[
 \begin{gathered}
 a+b+c=0,\qquad 0\le\ord(c)\le\ord(b)\le\ord(a)\le5,\\
 \ord(a)+\ord(b)+\ord(c)\le11.
 \end{gathered}
\]
Choose $a$ up to a coordinate permutation and sign by taking the lesser
of its sorted entries and the sorted entries of $-a$.  Enumerate all
balanced $b$ of orders at most $\ord(a)$, put $c=-a-b$, and retain those
meeting the displayed restrictions.  Sorting the six rows of the
array $(a,b,c)$ removes common-coordinate duplicates.  Balanced
vectors of a given order are generated by subtracting two
disjointly supported nonnegative compositions of that order, which is
the routine \texttt{vecs}.

For each generated array, \texttt{mixeq} first evaluates all five
equality patterns on its three labels, registering any single-label
zero equations.  It then considers the distinct permutations of the
three charge columns, choosing $a\le b$ lexicographically to avoid
duplicated first-two-column orderings.  Write
\[
 a_{ij}=a+e_i-e_j,\qquad b_{ij}=b-e_i+e_j
 \quad (0\le i,j\le5).
\]
\Needspace{6\baselineskip}
The following equations are included after applying \texttt{evmix} to
every term:
\begin{itemize}
 \item For different first two labels, sum the moments with charges
 $(a_{ij},b_{ij},c)$ over all $i\ne j$, obtaining zero by
 \eqref{eq:different-matrices}.  For nonzero $c$, use the label tuples
 $(0,1,2)$, $(0,1,0)$, and $(0,1,1)$; for $c=0$, use $(0,1,0)$.
 Include this equation only if all $36$ shifted triples have total
 order at most eleven.
 \item For equal first two labels, fix $i$, sum over all six $j$, and
 subtract the moment with charges $(-c,c)$ and labels $(0,u)$, by
 \eqref{eq:same-matrix}.  Use $(0,0,u)$ with $u=0,1$ when $c\ne0$,
 and $u=0$ when $c=0$.  Also use the reversed shifts
 $(a+e_j-e_i,b-e_j+e_i,c)$.  Include each equation only if all its six
 shifted triples have total order at most eleven.
\end{itemize}
The third charge is unchanged in these sums because the equations have
been multiplied by its factor.  On the right side of the equal-label
equation, $a+b=-c$, which explains both remaining charges and the
labels selected by \texttt{s[-2:]} in the source.

Finally add the pair equations recorded in \texttt{rows0}.  All
coordinates encountered by these reductions are registered, including
coordinates that cancel in a particular equation.  With the specified
coordinate order, \texttt{mixtable} applies
Lemma~\ref{lem:exact-elimination} to the nonzero equations.  A vanishing
expression encountered later may still be replaced by zero: its
vanishing is an independently proved rule, whether or not its pair
expansion was used to strengthen this table.  No completeness of the
enumeration, and no vanishing of unlisted higher-order moments, is
assumed.

The exact table dimensions are as follows.  The witness column counts
the independent equations used in Lemma~\ref{lem:exact-elimination}.
\begin{center}
\begin{tabular}{lrrr}
\toprule
Table & Nonzero equations & Coordinates & Witness rows \\
\midrule
Pair  & $8871$  & $909$  & $897$ \\
Mixed & $46041$ & $3429$ & $2767$ \\
\bottomrule
\end{tabular}
\end{center}

\subsection{Positive matrices and the negative bound}

For the mixed certificate, the literal data in Appendix~\ref{app:code}
specify the six lists and lower triangles of the symmetric integer
matrices $Z_\lambda$.  The routine \texttt{makeZ} checks the number of
entries and reconstructs each symmetric matrix.  Positive definiteness
is verified without eigenvalue approximation.  At each step,
\texttt{posdef} writes the remaining matrix as
\[
 B=\begin{pmatrix}C&u\\u^{\mathsf T}&p\end{pmatrix},
 \qquad p>40000,
\]
and replaces it by the exact Schur complement $C-uu^{\mathsf T}/p$.
The identity
\[
 \begin{pmatrix}x\\t\end{pmatrix}^{\!*}
 B\begin{pmatrix}x\\t\end{pmatrix}
 =x^*(C-uu^{\mathsf T}/p)x
   +p\left|t+\frac{u^{\mathsf T}x}{p}\right|^2
\]
shows inductively that positive pivots prove positive definiteness,
also for complex vectors.  All $98$ pivots in the six matrices satisfy
the asserted rational inequality.

The routines \texttt{permsign} and \texttt{polynomial} implement the
signed permutation sums defining $Y_\lambda$ in
Section~\ref{sec:mixed}.  The convention
$(P_pv)_i=v_{p(i)}$ gives
$(P_pP_qv)_i=v_{q(p(i))}$, exactly
\texttt{perm(perm(v,q),p)}.  The routine \texttt{tagmom} uses the Gram
charges $(a,-v,v-a)$ of Lemma~\ref{lem:gram}.  After multiplying the
entire form by five, \texttt{product} assigns weight five unless both
function types are one; in that case the weights for equal and
different external labels are one and four.  The lower-triangle
summation in \texttt{gaps} doubles the off-diagonal terms.  Since the
polynomial coefficients and $Z_\lambda$ are real, its real-moment
expansion equals the full quadratic form for complex $Y_\lambda$.

Consequently the expression formed by \texttt{gaps} is exactly
\[
 \mathcal Q=
 5\sum_\lambda 10^{15-e_\lambda}
       \E\bigl(Y_\lambda^*Z_\lambda Y_\lambda\bigr).
\]
Apply \texttt{evmix} and then the certified mixed substitutions.  On
combining coefficients, we obtain the residual expansion in
\eqref{eq:mixed-residual}, now indexed by the program's moment keys:
\begin{equation}
 \begin{gathered}
 \mathcal Q=c_0+\sum_{K\ne\varnothing}c_K\,\mu_K,\\
 c_0=-\frac{613302797399911}{6480},\qquad
 \sum_{K\ne\varnothing}|c_K|
     =\frac{179721388988719}{6480}.
 \end{gathered}
 \label{eq:exact-residual-coefficients}
\end{equation}
Here $\mu_K$ denotes the original moment named by a retained coordinate;
the notation $K\ne\varnothing$ excludes the completely empty constant
key.  Both equalities of rational coefficients in
\eqref{eq:exact-residual-coefficients} are checked exactly by the
program, as is their strictly negative sum.

It remains to justify the bound on these residual coordinates.  A
tagged pair coordinate is a real average of unit-modulus monomials, so
$|m(A)|\le1$.  A mixed coordinate is a real average of products of
$g_s(a)$, each of absolute value at most one, so
$|\mathcal M(\mathbf s;\mathbf a)|\le1$.  Every coordinate retained by
the substitutions is one of these original moments.  The new linear
combinations occur as expressions in the coordinates, not as newly
named coordinates to which a bound is assigned.  The constant
coordinate is exactly one.  Hence
\[
 0\le\mathcal Q
 \le c_0+\sum_{K\ne\varnothing}|c_K|
 =-\frac{2007321335237}{30}<0.
\]
This proves Proposition~\ref{prop:mixed-certificate}.  Together with
the pair coefficient check, it completes the exact verification of
both certificates using only the finite equations and integer data
specified in Appendix~\ref{app:code}.
